\chapter{The Problem and Its Two Observation Regimes}
\label{ch:problem}

\section{What we are asked to predict}

An investment platform tracking private companies wants, at any date $t$, answers to three
nested questions:
(i)~\emph{which private companies are succeeding}, in the sense of growing fundamentals
(EBITDA, revenue, employment) or achieving valuation-relevant milestones;
(ii)~\emph{which companies are likely to transact} --- raise a venture round, be acquired,
take on debt --- over a horizon $(t, t+h]$; and
(iii)~\emph{how market conditions} (rates, inflation, unemployment, risk appetite)
\emph{modulate both}. The three questions look like one prediction problem, but they are
not: what separates them is the \emph{observation process}, i.e.\ which state variables of
the firm the data vendor lets us see and \emph{when}.

This manual is organised around that separation. We work with two data assets. The first
is a deals-and-companies database in the PitchBook schema: a \texttt{deals} table in which
each row is a transaction (with a \texttt{deal\_class} taxonomy --- venture capital rounds,
M\&A/buyouts, debt financings --- each populating a different subset of fields), a
\texttt{companies} table describing the tracked universe, and an \texttt{investors}
relation. The second is a standardized macroeconomic feed in the FactSet Standardized
Economics convention, from which we use the policy-rate series (\texttt{FFUNDT},
\texttt{FFUND}), consumer prices (\texttt{CPI\_RAW}, \texttt{CPI\_YOY}) and unemployment
(\texttt{RUNEMP}), each published with a realistic lag. Chapter~\ref{ch:data} describes
both schemas in detail, together with the point-in-time discipline that governs every join
in this manual.

\section{Two regulatory regimes, two modelling worlds}

\paragraph{Europe: fundamentals are observable.}
In most European jurisdictions, limited-liability companies must file annual accounts ---
balance sheet, income statement, often cash-flow --- with a public registry, essentially
irrespective of size. For a European private company we therefore observe a
\emph{fundamentals panel}: $F_{i,y} = (\text{revenue}, \text{EBITDA}, \text{payroll},
\text{assets}, \dots)$ at yearly frequency, with publication lags but near-complete
coverage. ``Success'' can be defined directly on this panel (Chapter~\ref{ch:success}),
predicted with supervised learning on engineered features (Chapter~\ref{ch:xgb}), and
interrogated causally --- \emph{does raising a round cause EBITDA growth, or do growing
firms raise?} --- with the backdoor criterion and bootstrap uncertainty quantification
(Chapter~\ref{ch:causal}).

\paragraph{United States: fundamentals are latent.}
US private companies file almost nothing publicly. Revenue or EBITDA appear in the deals
data only sporadically (in our synthetic mirror, calibrated to vendor coverage,
roughly 45\% of deal rows carry a revenue figure and fewer an EBITDA). What we observe
instead is the \emph{transaction record}: when a firm raised, at what valuation, which
round, how many employees it reported at that moment. The firm's state is measured
\emph{only at its own deal events}; between events we can only carry the last observation
forward. Part~II develops what can still be said about growth in this regime: deal marks
as noisy state proxies, the decomposition of valuation growth into fundamentals and market
multiple, and macro factors as partially exogenous drivers.

\paragraph{The event-based endgame.}
Once one accepts that the US data are a \emph{marked event stream} --- timestamps plus
deal-level marks, refreshed only at events --- the natural formulation is no longer
regression on a panel but \emph{calibration of a counting process}: an intensity
$\lambda(t)$ for deal arrivals, modulated by covariates, history (excitation and
inhibition), and firm identity. Part~III builds this up in deliberately increasing order
of structure: inhomogeneous Poisson (Chapter~\ref{ch:poisson}), Hawkes processes with
excitation and the two-stage marked decomposition (Chapter~\ref{ch:hawkes}), and
interval-censored variants for when even the timestamps are unreliable
(Chapter~\ref{ch:censored}). Chapter~\ref{ch:selection} closes with decision tables and a
synthetic-testbed validation of every major claim.

\section{Guiding principles}

Four principles recur in every chapter; we state them once here.

\begin{enumerate}
\item \textbf{Point-in-time discipline.} Every feature used at time $t$ must be computable
from data published strictly before $t$. Macro series join on \emph{publish date}, not
reference date; firm features are last-observation-carried-forward from \emph{past} deals;
end-of-sample snapshot columns (e.g.\ \texttt{last\_financing\_size} in the companies
table) are forbidden as features. Violations produce leakage that flatters every offline
metric and survives until production.
\item \textbf{Identifiability before estimation.} We ask \emph{which} parameter
combinations the data can pin down before asking how to estimate them. Recurring examples:
covariate effects are identified by contrast while baseline levels are weak; branching
ratios are identified from counts while decay rates need within-interval timing;
excitation is confounded by latent common factors. Each claim is made precise where it
arises and verified on synthetic ground truth.
\item \textbf{Uncertainty quantification is part of the deliverable.} Point predictions
without calibrated uncertainty are not usable for ranking or allocation. We use, by
context: sandwich/quasi-likelihood corrections for over-dispersed counts, cluster
bootstrap for firm-level dependence, quantile/conformal methods for tabular prediction,
and Laplace/posterior summaries for point-process parameters.
\item \textbf{Causal claims are labelled as such.} Predictive and causal targets are kept
notationally distinct ($\E[Y \mid X{=}x]$ vs.\ $\E[Y \mid \Do(X{=}x)]$). Where a causal
reading is wanted --- the effect of financing on growth, of rates on deal flow, of past
deals on future deals (``contagion'') --- we state the identification assumptions, test
what is testable, and demonstrate the full pipeline on synthetic data where the true
effect is known.
\end{enumerate}

\section{Notation}
\label{sec:notation}

\begin{table}[htbp]
\centering
\small
\begin{tabular}{ll}
\toprule
Symbol & Meaning \\
\midrule
$i = 1,\dots,N$ & firms; $s(i) \in \{1,\dots,M\}$ the sector of firm $i$ \\
$t$ & calendar time (days unless stated); $w$ indexes weeks, $y$ years \\
$F_{i,y}$ & fundamentals vector of firm $i$ in fiscal year $y$ (EU regime) \\
$Y$, $A$, $X$, $U$ & outcome, treatment/exposure, observed covariates, unobserved factors \\
$N_i(t)$ & counting process of firm $i$'s deals; $N_s(t)$ sector aggregate \\
$\Hist_t$ & information (filtration) available just before $t$ \\
$\lambda_i(t)$ & conditional intensity of $N_i$ given $\Hist_t$ \\
$\Lambda(a,b] = \int_a^b \lambda(u)\dd u$ & compensator increment \\
$s(t)$, $\phi(\cdot)$ & baseline intensity; excitation kernel \\
$\kappa$, $\theta$ & branching ratio and decay rate of an exponential kernel \\
$G = \sum_\ell b_\ell$, $\rho(G)$ & aggregate excitation matrix and spectral radius \\
$Z_{i,t}$ & firm-level point-in-time features; $X_t$ market covariates \\
$\RS_{s,t}$ & risk set (candidate firms) for a sector-$s$ event at $t$ \\
$q_{i,t}$ & choice-model score of candidate $i$; $\eta$ cooldown/inhibition coefficient \\
$\E[\cdot], \Var, \Prob$ & expectation, variance, probability; $\ind{\cdot}$ indicator \\
$\Do(\cdot)$ & Pearl's intervention operator \\
\bottomrule
\end{tabular}
\caption{Global notation. Chapter-local symbols are defined where used.}
\label{tab:notation}
\end{table}

\section{The synthetic testbed}
\label{sec:testbed}

Every method in this manual is exercised on a synthetic world that mirrors the production
schemas exactly (all PitchBook deal fields, the FactSet long-format feed with publication
lags) and whose ground truth we control. Two regimes are shipped with the code:
\textbf{regime~A}, generated by the exact two-stage weekly model of
Chapter~\ref{ch:hawkes} (a well-specified world for parameter-recovery studies), and
\textbf{regime~B}, generated by a daily firm-level process with self-inhibition (30-day
half-life), sector contagion (7-day half-life) and latent firm quality (a misspecified
world for robustness studies). Both contain roughly $50{,}000$ deals over ten years,
twelve sectors and $8{,}000$ firms, with realistic observation noise: 28\% undisclosed
deal sizes, 55\% missing valuations, 3\% mislabelled sectors, duplicated records, jittered
dates, and a 12\% untracked (out-of-universe) firm share. The accompanying
\texttt{code/} directory contains the generator, point-in-time loaders, and fast reference
fitters; Chapter~\ref{ch:selection} tabulates the measured results that are quoted
throughout the text.

\begin{remark}
Numbers quoted from the testbed (e.g.\ recovery correlations, held-out likelihoods) are
single-seed values from the shipped datasets unless stated otherwise; the experiment
protocol in the code repeats them over five seeds.
\end{remark}
